% Distilled mathematical notes for:
%
%   Yi Yu, Tengyao Wang, and Richard J. Samworth,
%   "A useful variant of the Davis--Kahan theorem for statisticians",
%   arXiv:1405.0680 [math.ST], submitted 4 May 2014.
%   DOI: 10.48550/arXiv.1405.0680
%
% This file is not a transcription of the article.  It is an independent
% distillation of the core mathematical argument, with attribution.
%
% Notation switch:
%   \\PaperNotationtrue  : Yu--Wang--Samworth notation, e.g. Sigma, hat Sigma.
%   \\PaperNotationfalse : generic perturbation notation, e.g. A, tilde A.

\documentclass[11pt]{article}

\usepackage[T1]{fontenc}
\usepackage{lmodern}
\usepackage{microtype}
\usepackage[margin=1.1in]{geometry}
\usepackage{amsmath,amssymb,amsthm,mathtools}
\usepackage[hidelinks]{hyperref}
\usepackage{enumitem}

% -----------------------------------------------------------------------------
% Notation layer
% -----------------------------------------------------------------------------
\newif\ifPaperNotation
\PaperNotationtrue

\ifPaperNotation
  \newcommand{\BaseMat}{\Sigma}
  \newcommand{\PertMat}{\widehat{\Sigma}}
  \newcommand{\ErrMat}{\widehat{\Sigma}-\Sigma}
  \newcommand{\Sub}{V}
  \newcommand{\Subhat}{\widehat V}
  \newcommand{\EigDiag}{\Lambda}
  \newcommand{\EigDiagHat}{\widehat\Lambda}
  \newcommand{\AlignMat}{\widehat O}
\else
  \newcommand{\BaseMat}{A}
  \newcommand{\PertMat}{\widetilde A}
  \newcommand{\ErrMat}{\widetilde A-A}
  \newcommand{\Sub}{Q}
  \newcommand{\Subhat}{\widetilde Q}
  \newcommand{\EigDiag}{D}
  \newcommand{\EigDiagHat}{\widetilde D}
  \newcommand{\AlignMat}{O}
\fi

\newcommand{\R}{\mathbb R}
\newcommand{\gap}{\Delta}
\newcommand{\sgap}{\Gamma}
\newcommand{\norm}[1]{\left\lVert #1 \right\rVert}
\newcommand{\op}{\mathrm{op}}
\newcommand{\F}{\mathrm F}
\newcommand{\diag}{\operatorname{diag}}
\newcommand{\tr}{\operatorname{tr}}
\newcommand{\vecop}{\operatorname{vec}}
\newcommand{\rank}{\operatorname{rank}}

\newtheorem{theorem}{Theorem}
\newtheorem{lemma}[theorem]{Lemma}
\newtheorem{corollary}[theorem]{Corollary}
\theoremstyle{remark}
\newtheorem*{remark}{Remark}

\title{Core arguments in Yu--Wang--Samworth's Davis--Kahan variant}
\author{Distilled notes}
\date{}

\begin{document}
\maketitle

\begin{abstract}
These notes isolate the main mathematical mechanism in Yu, Wang, and
Samworth's useful Davis--Kahan variant.  The key observation is that the
residual $\Subhat\EigDiag-\BaseMat\Subhat$ can be bounded above using only
matrix perturbation estimates, while it can be bounded below by a
population eigengap times the sine of the principal angles.  The result is
an eigenvector/subspace perturbation bound whose denominator contains only
population eigenvalue gaps.
\end{abstract}

\section{Reference and scope}

The source paper is:
\begin{quote}
Yi Yu, Tengyao Wang, and Richard J. Samworth,
``A useful variant of the Davis--Kahan theorem for statisticians,''
arXiv:1405.0680, 2014.
\end{quote}

This document is a short mathematical distillation, not a line-by-line
transcription.  It keeps the paper's notation by default.  To view the same
argument in more generic perturbation notation, set
\verb|\PaperNotationfalse| near the top of the file.

\section{Setup}

Let $\BaseMat,\PertMat\in\R^{p\times p}$ be symmetric, with eigenvalues
$\lambda_1\geq\cdots\geq\lambda_p$ and
$\widehat\lambda_1\geq\cdots\geq\widehat\lambda_p$, respectively.  Fix
indices $1\leq r\leq s\leq p$ and write $d=s-r+1$.  The target population
subspace is spanned by
\[
  \Sub=(v_r,v_{r+1},\ldots,v_s)\in\R^{p\times d},
\]
and the perturbed subspace is spanned by
\[
  \Subhat=(\widehat v_r,\widehat v_{r+1},\ldots,\widehat v_s)
  \in\R^{p\times d}.
\]
Thus
\[
  \BaseMat\Sub=\Sub\EigDiag,
  \qquad
  \PertMat\Subhat=\Subhat\EigDiagHat,
\]
where
\[
  \EigDiag=\diag(\lambda_r,\ldots,\lambda_s),
  \qquad
  \EigDiagHat=\diag(\widehat\lambda_r,\ldots,\widehat\lambda_s).
\]
The relevant population eigengap is
\[
  \gap := \min(\lambda_{r-1}-\lambda_r,\lambda_s-\lambda_{s+1}),
  \qquad
  \lambda_0:=\infty,\quad \lambda_{p+1}:=-\infty.
\]
The goal is to control the principal angles between the column spaces of
$\Sub$ and $\Subhat$, measured by $\norm{\sin\Theta(\Subhat,\Sub)}_{\F}$.

\section{The symmetric-matrix variant}

\begin{theorem}[Population-gap Davis--Kahan variant]
Assume $\gap>0$.  Then
\[
  \norm{\sin\Theta(\Subhat,\Sub)}_{\F}
  \leq
  \frac{2\min\{d^{1/2}\norm{\ErrMat}_{\op},\norm{\ErrMat}_{\F}\}}{\gap}.
\]
Moreover, there is a $d\times d$ orthogonal matrix $\AlignMat$ such that
\[
  \norm{\Subhat\AlignMat-\Sub}_{\F}
  \leq
  \frac{2^{3/2}\min\{d^{1/2}\norm{\ErrMat}_{\op},\norm{\ErrMat}_{\F}\}}{\gap}.
\]
\end{theorem}

\begin{remark}
The important structural point is that $\gap$ is formed only from the
population eigenvalues.  The usual statistical workaround is to apply a
classical Davis--Kahan theorem with a mixed population/sample separation,
and then use Weyl's inequality to show that the sample eigenvalues are close
to the population eigenvalues.  The theorem above builds the population gap
directly into the argument.
\end{remark}

\section{The proof as an upper/lower sandwich}

Define the residual
\[
  R := \Subhat\EigDiag-\BaseMat\Subhat.
\]
The whole proof is the comparison
\[
  \gap\,\norm{\sin\Theta(\Subhat,\Sub)}_{\F}
  \leq \norm{R}_{\F}
  \leq 2\min\{d^{1/2}\norm{\ErrMat}_{\op},\norm{\ErrMat}_{\F}\}.
\]
The left inequality is a deterministic separation estimate.  The right
inequality is a perturbation estimate.

\subsection{Upper bound on the residual}

From $\PertMat\Subhat=\Subhat\EigDiagHat$ we have
\[
  0=\PertMat\Subhat-\Subhat\EigDiagHat
   =\BaseMat\Subhat-\Subhat\EigDiag+(\ErrMat)\Subhat
     -\Subhat(\EigDiagHat-\EigDiag),
\]
and hence
\[
  R=(\ErrMat)\Subhat-\Subhat(\EigDiagHat-\EigDiag).
\]
Therefore
\[
  \norm{R}_{\F}
  \leq \norm{(\ErrMat)\Subhat}_{\F}
      +\norm{\Subhat(\EigDiagHat-\EigDiag)}_{\F}.
\]
There are two useful versions of this upper bound.

First, using $\norm{B\Subhat}_{\F}\leq d^{1/2}\norm{B}_{\op}$ and Weyl's
inequality,
\[
  \norm{R}_{\F}
  \leq d^{1/2}\norm{\ErrMat}_{\op}+\norm{\EigDiagHat-\EigDiag}_{\F}
  \leq 2d^{1/2}\norm{\ErrMat}_{\op}.
\]
Second, using Frobenius-norm contraction under multiplication by matrices
with orthonormal columns, together with the Hoffman--Wielandt eigenvalue
perturbation bound,
\[
  \norm{R}_{\F}
  \leq \norm{\ErrMat}_{\F}+\norm{\EigDiagHat-\EigDiag}_{\F}
  \leq 2\norm{\ErrMat}_{\F}.
\]
Combining the two gives
\[
  \norm{R}_{\F}
  \leq 2\min\{d^{1/2}\norm{\ErrMat}_{\op},\norm{\ErrMat}_{\F}\}.
\]

\subsection{Lower bound on the residual}

Let $V_1\in\R^{p\times(p-d)}$ be an orthonormal complement to $\Sub$ chosen
so that, with $P=(\Sub\; V_1)$,
\[
  P^T\BaseMat P=\begin{pmatrix}\EigDiag&0\\0&\Lambda_1\end{pmatrix},
  \qquad
  \Lambda_1=\diag(\lambda_1,\ldots,\lambda_{r-1},
                   \lambda_{s+1},\ldots,\lambda_p).
\]
Projecting $R$ onto the orthogonal complement of $\operatorname{col}(\Sub)$
can only decrease its Frobenius norm, so
\[
  \norm{R}_{\F}
  \geq
  \norm{V_1^T\Subhat\EigDiag-\Lambda_1V_1^T\Subhat}_{\F}.
\]
The linear map
\[
  X\mapsto X\EigDiag-\Lambda_1X
\]
acts entrywise by multiplying the $(a,b)$ entry of $X$ by a difference
between one eigenvalue outside the target block and one eigenvalue inside
it.  By the definition of $\gap$, all these multipliers have absolute value
at least $\gap$.  Hence
\[
  \norm{V_1^T\Subhat\EigDiag-\Lambda_1V_1^T\Subhat}_{\F}
  \geq \gap\,\norm{V_1^T\Subhat}_{\F}.
\]
Finally,
\[
  \norm{V_1^T\Subhat}_{\F}^2
  =\tr\{\Subhat^T(I_p-\Sub\Sub^T)\Subhat\}
  =d-\norm{\Subhat^T\Sub}_{\F}^2
  =\norm{\sin\Theta(\Subhat,\Sub)}_{\F}^2.
\]
Thus
\[
  \norm{R}_{\F}
  \geq \gap\,\norm{\sin\Theta(\Subhat,\Sub)}_{\F}.
\]

\section{From angles to aligned bases}

The sine-angle bound controls the subspaces.  To control a particular choice
of orthonormal bases, align them by an orthogonal Procrustes rotation.  Let
\[
  \widehat O_1^T\Subhat^T\Sub\widehat O_2
  =\diag(\cos\theta_1,\ldots,\cos\theta_d)
\]
be a singular value decomposition, and set
$\AlignMat=\widehat O_1\widehat O_2^T$.  Then
\[
\begin{aligned}
  \norm{\Subhat\AlignMat-\Sub}_{\F}^2
  &=2d-2\sum_{j=1}^d\cos\theta_j  \\
  &\leq 2d-2\sum_{j=1}^d\cos^2\theta_j
   =2\norm{\sin\Theta(\Subhat,\Sub)}_{\F}^2.
\end{aligned}
\]
Taking square roots and inserting the previous theorem gives the displayed
basis-alignment bound.

\section{The single-vector corollary}

When $d=1$, say the target index is $j$, the theorem reduces to
\[
  \sin\Theta(\widehat v_j,v_j)
  \leq
  \frac{2\norm{\ErrMat}_{\op}}
       {\min(\lambda_{j-1}-\lambda_j,\lambda_j-\lambda_{j+1})}.
\]
If the sign of $\widehat v_j$ is chosen so that $\widehat v_j^Tv_j\geq 0$,
then
\[
  \norm{\widehat v_j-v_j}
  \leq \sqrt2\,\sin\Theta(\widehat v_j,v_j)
  \leq
  \frac{2^{3/2}\norm{\ErrMat}_{\op}}
       {\min(\lambda_{j-1}-\lambda_j,\lambda_j-\lambda_{j+1})}.
\]

\section{The singular-vector extension}

Let $A,\widehat A\in\R^{p\times q}$ have singular values
$\sigma_1\geq\cdots\geq\sigma_{\min(p,q)}$ and
$\widehat\sigma_1\geq\cdots\geq\widehat\sigma_{\min(p,q)}$.  Fix
$1\leq r\leq s\leq\rank(A)$, set $d=s-r+1$, and define the squared-singular
value gap
\[
  \sgap := \min(\sigma_{r-1}^2-\sigma_r^2,
                 \sigma_s^2-\sigma_{s+1}^2),
  \qquad
  \sigma_0^2:=\infty,
  \quad
  \sigma_{\rank(A)+1}^2:=-\infty.
\]
Let $V=(v_r,\ldots,v_s)$ and $\widehat V=(\widehat v_r,\ldots,\widehat v_s)$
be the corresponding right singular-vector matrices.  Applying the symmetric
result to $A^TA$ and $\widehat A^T\widehat A$ gives
\[
  \norm{\sin\Theta(\widehat V,V)}_{\F}
  \leq
  \frac{2\min\{d^{1/2}\norm{\widehat A^T\widehat A-A^TA}_{\op},
               \norm{\widehat A^T\widehat A-A^TA}_{\F}\}}{\sgap}.
\]
The Gram-matrix perturbation is bounded by
\[
\begin{aligned}
  \norm{\widehat A^T\widehat A-A^TA}_{\op}
  &\leq (2\sigma_1+\norm{\widehat A-A}_{\op})\norm{\widehat A-A}_{\op},\\
  \norm{\widehat A^T\widehat A-A^TA}_{\F}
  &\leq (2\sigma_1+\norm{\widehat A-A}_{\op})\norm{\widehat A-A}_{\F}.
\end{aligned}
\]
Consequently,
\[
  \norm{\sin\Theta(\widehat V,V)}_{\F}
  \leq
  \frac{2(2\sigma_1+\norm{\widehat A-A}_{\op})
        \min\{d^{1/2}\norm{\widehat A-A}_{\op},
              \norm{\widehat A-A}_{\F}\}}{\sgap}.
\]
The same Procrustes step yields an orthogonal $\widehat O\in\R^{d\times d}$
for which
\[
  \norm{\widehat V\widehat O-V}_{\F}
  \leq
  \frac{2^{3/2}(2\sigma_1+\norm{\widehat A-A}_{\op})
        \min\{d^{1/2}\norm{\widehat A-A}_{\op},
              \norm{\widehat A-A}_{\F}\}}{\sgap}.
\]
Identical bounds for left singular subspaces follow by applying the same
argument to $AA^T$ and $\widehat A\widehat A^T$.

\section{Minimal supporting lemmas}

\begin{lemma}[Frobenius compression]
Let $B\in\R^{m\times n}$.  If $U\in\R^{m\times a}$ and
$W\in\R^{n\times b}$ have orthonormal columns, then
\[
  \norm{U^TBW}_{\F}\leq \norm{B}_{\F}.
\]
If instead $U$ and $W$ have orthonormal rows of the compatible sizes, then
\[
  \norm{U^TBW}_{\F}=\norm{B}_{\F}.
\]
\end{lemma}

\begin{proof}
Complete the column-orthonormal matrices to orthogonal matrices and observe
that Frobenius norm is preserved by orthogonal multiplication; the block
$U^TBW$ is one sub-block of the transformed matrix, hence has no larger
Frobenius norm.  For orthonormal rows, $UU^T=I$ and $WW^T=I$, so
\[
  \norm{U^TBW}_{\F}^2
  =\tr(U^TBWW^TB^TU)=\tr(BB^TUU^T)=\tr(BB^T)=\norm{B}_{\F}^2.
\]
\end{proof}

\begin{lemma}[Eigenvalue perturbation inputs]
For symmetric matrices $M$ and $N$ with ordered eigenvalues $\mu_i$ and
$\nu_i$,
\[
  \max_i |\mu_i-\nu_i|\leq \norm{M-N}_{\op},
  \qquad
  \sum_i(\mu_i-\nu_i)^2\leq \norm{M-N}_{\F}^2.
\]
The first inequality is Weyl's inequality; the second is the
Hoffman--Wielandt inequality in the symmetric case.
\end{lemma}

\section{Takeaway}

The proof is not a contour-integral or resolvent argument.  It is a residual
argument:
\[
  \text{population separation} \times \text{subspace error}
  \leq \text{residual}
  \leq \text{matrix perturbation size}.
\]
That is why the denominator can be a population gap only.  The price is a
small constant factor, while the benefit in statistical applications is that
one can state assumptions directly on the population spectrum.

\begin{thebibliography}{9}
\bibitem{YWS2014}
Yi Yu, Tengyao Wang, and Richard J. Samworth.
\newblock A useful variant of the Davis--Kahan theorem for statisticians.
\newblock \emph{arXiv:1405.0680} [math.ST], 2014.
\newblock \url{https://arxiv.org/abs/1405.0680}.

\bibitem{DavisKahan1970}
Chandler Davis and W. M. Kahan.
\newblock The rotation of eigenvectors by a perturbation. III.
\newblock \emph{SIAM Journal on Numerical Analysis}, 7(1):1--46, 1970.
\end{thebibliography}

\end{document}
